% Chapter 2: Data structure and algorithm specification.

\section{Data Structure and Algorithm Specification}
\label{sec:design}

This chapter specifies the data structures and algorithms used by the engine.
The order follows the data flow: tokenization, counting, stop-word selection,
stemming, index construction, dictionary structures, disk-based storage,
external construction, and query processing. Pseudo-code is written in a
language-independent form; identifiers such as \code{lastDoc} and
\code{posStart} match the implementation.

\subsection{Tokenization and Normalization}
\label{sec:tokenization}

The corpus is stored as a \emph{text layer}: one semantic unit per line, with a
line marker that identifies headings, stage directions, and speaker labels.
\code{manifest.tsv} is the catalog; each row contains the document ID, title,
text path, and source path, and an optional \code{format} column selects
\code{marked} (default) or \code{plain} (used by the RFC corpus).

A token is a maximal run of ASCII letters that may contain internal
apostrophes. All tokens are lowercased. The apostrophe rules are:

\begin{itemize}
  \item an internal apostrophe is kept: \code{don't} and \code{o'er} are single
        tokens;
  \item a leading apostrophe is dropped: \code{'Tis} becomes \code{tis};
  \item a trailing apostrophe acts as a separator: \code{friends'} becomes
        \code{friends};
  \item hyphens, digits, and ordinary punctuation are separators:
        \code{to-day} becomes \code{to} and \code{day}, and \code{1660s}
        becomes \code{s}.
\end{itemize}

Marker lines are skipped by default. A heading such as \code{[ACT I]}, a stage
direction such as \code{(Exit Ghost)}, or a speaker label such as
\code{@BERNARDO} is recognized only when the first and last characters match
the expected pair. A line that begins with \code{[} but continues as dialogue
is body text. The optional flags \code{--keep-headings}, \code{--keep-speakers},
and \code{--keep-stage} retain each class for experiments that need a different
counting convention.

\begin{lstlisting}[style=cli]
for each line L in the text layer:
    if L is a marker line and the marker class is not kept:
        continue
    for each maximal run T of ASCII letters and internal apostrophes:
        T = drop leading apostrophe; T = lowercase(T)
        emit T
\end{lstlisting}

Normalization is the composition of three stages: tokenize, optionally remove
stop words, and optionally apply the Porter stemmer. The index header stores
two flags for the latter two stages, and the query program reads these flags
and replays the same composition. This is the mechanism that keeps a query
from using a different definition of a term than the index.

\subsection{Word Frequency Statistics}
\label{sec:count-design}

\subsubsection{Design Principles}

For a term $t$, the engine reports its term frequency and document frequency:
\begin{equation}
  \mathrm{tf}(t)=\sum_{d}\mathrm{tf}(t,d),\qquad
  \mathrm{df}(t)=|\{d:\mathrm{tf}(t,d)>0\}|.
\end{equation}
It also computes the document-frequency ratio
$\mathrm{df\_ratio}(t)=\mathrm{df}(t)/195$ and the term-frequency share
$\mathrm{tf\_share}(t)=\mathrm{tf}(t)/866\,626$. The first measures coverage;
the second measures volume. Looking at only one of them makes a dialogue word
diluted by the 154 short poems look rare when it is actually common.

The counter reuses the formal normalization pipeline. Four configurations are
supported: the default (no stop words, no stemming), stop words only, stemming
only, and the formal configuration (both). The same order---tokenize, filter,
stem---is used by the index and query programs.

\subsubsection{Algorithm}

Documents are scanned in manifest order. Each term row stores its accumulated
frequency and the ID of the last document in which it was seen. The document
frequency is incremented only when the current document ID differs from
\code{lastDoc}; repeated occurrences inside one document therefore contribute
once to df.

\begin{lstlisting}[style=cli]
for each document d in manifest order:
    for each token t in tokenize(d):
        row = table[t]
        row.tf += 1
        if row.lastDoc != d.id:
            row.df += 1
            row.lastDoc = d.id
\end{lstlisting}

With a hash table the expected time is $O(N)$ for $N$ tokens, plus $O(T\log T)$
for sorting the output rows; with a splay dictionary it is $O(N\log T)$
amortized. The memory bound is $O(T)$ rows. The four CSV outputs are
deterministic: the same input and flags produce the same bytes, and the files
contain no timestamps.

\subsection{Stop-Word Selection}
\label{sec:stopwords}

\subsubsection{The tf/df Candidate Rule}

The simplest candidate rule asks a word to be both widespread and frequent:
\begin{equation}
  \mathrm{df\_ratio}(t)\ge 0.75
  \quad\land\quad
  \mathrm{tf\_share}(t)\ge 0.001.
\end{equation}
It returns only nine words. The reason is the corpus composition: 154 short
poems account for 79\% of the documents, and many function words do not appear
in every poem, so their document-frequency ratios are diluted. For example,
\code{a} has a ratio of 0.7077, just below 0.75. Lowering the threshold to 0.40
produces 48 candidates, but content words such as \code{love} and \code{time}
enter the band as well.

\subsubsection{Length-Normalized Rate Entropy}
\label{sec:entropy}

To separate ``spread evenly across many documents'' from ``merely frequent,''
the engine computes a length-normalized rate entropy. Let $L_d$ be the number
of kept tokens in document $d$:
\begin{equation}
  r_{t,d}=\frac{\mathrm{tf}(t,d)}{L_d},\qquad
  q_{t,d}=\frac{r_{t,d}}{\sum_{d'}r_{t,d'}},\qquad
  \widehat H(t)=\frac{-\sum_{d}q_{t,d}\log_2 q_{t,d}}{\log_2\mathrm{df}(t)}.
  \label{eq:rate-entropy}
\end{equation}
After length normalization, function words sit above 0.95, content words such
as \code{love}, \code{king}, and \code{heart} are visibly lower, and a proper
noun with $\mathrm{df}=1$ has entropy zero by definition. The denominator
$\log_2\mathrm{df}(t)$ makes the statistic comparable across words that occur
in different numbers of documents.

\subsubsection{The Three-Feature Rule}

The final rule requires a word to be wide, heavy, and flat:
\begin{equation}
  \mathrm{df\_ratio}\ge r
  \quad\land\quad
  \mathrm{tf\_share}\ge s
  \quad\land\quad
  \widehat H\ge h.
\end{equation}
The default $(r,s,h)=(0.75,0.001,0.93)$ yields nine words. The experimental
configuration $(0.40,0.001,0.93)$ yields 39 words and overlaps the manually
curated list in 35 words. The automatic list adds \code{hath}, \code{o},
\code{thee}, and \code{thou}, and omits \code{more}, \code{or}, \code{than},
and \code{their}; each difference can be explained by the rate entropy.

\begin{lstlisting}[style=cli]
for each term t in the vocabulary:
    if df_ratio[t] >= r and tf_share[t] >= s
       and rate_entropy[t] >= h:
        emit t
\end{lstlisting}

The rule is deterministic and uses no external stop-word list. Changing a
threshold changes only the strictness of the band, not the structure of the
rule.

\subsection{Porter Stemming}
\label{sec:stemmer}

\subsubsection{Comparison of Public Stemmers}

Four public descriptions were compared: the original Porter (1980) paper, the
canonical C implementation on the author's site, the NLTK
\code{PorterStemmer} in its three compatibility modes, and the Snowball Porter2
algorithm \cite{porter1980,porter-c,nltk,snowball}. The implementation follows
the 1980 paper plus the three marked \code{--DEPARTURE--} changes in the
canonical C code: words of two letters or fewer are left unchanged, \code{bli}
is changed to \code{ble}, and \code{logi} is changed to \code{log}. This is the
same behavior as NLTK's \code{MARTIN_EXTENSIONS} mode, which the full-vocabulary
cross-check uses as an oracle.

\subsubsection{Algorithm}

Before the Porter rules, contractions are reduced: the longest matching suffix
among \code{n't}, \code{'s}, \code{'d}, \code{'ll}, \code{'re}, \code{'ve},
\code{'t}, and \code{'m} is removed, and remaining apostrophes are deleted.
Thus \code{o'er} becomes \code{oer}, \code{marry'd} becomes \code{marryd}, and
the irregular forms keep their rule result, for example \code{can't} becomes
\code{ca} and \code{won't} becomes \code{wo}.

The main flow executes steps 1a, 1b, 1c, 2, 3, 4, 5a, and 5b. The core
condition is the measure $m$ (the number of vowel-consonant sequences) and the
predicates \code{*v*}, \code{*d}, and \code{*o}. The implementation follows
``the first matching suffix decides'': once a suffix matches, only that rule's
condition is checked; if the condition fails, the step stops and does not fall
back to a shorter suffix. This detail matters: an earlier version continued to
shorter suffixes after a failed condition, over-stemmed 41 word forms (for
example, \code{agreement} became \code{agreem}), and was caught by the
full-vocabulary differential test.

\begin{lstlisting}[style=cli]
stem(word):
    word = strip_longest_contraction(word)
    for step in (1a, 1b, 1c, 2, 3, 4, 5a, 5b):
        for suffix in step.suffixes_by_length:
            if word ends with suffix:
                if step.condition(word, suffix):
                    word = step.replace(word, suffix)
                break          # first match wins; do not try shorter suffixes
    return word
\end{lstlisting}

The worst-case cost is $O(L^2)$ for a word of length $L$ because suffix
matching and condition checks rescan the word; in practice $L$ is small and the
observed throughput is dominated by the dictionary, not by the stemmer.

\subsection{Inverted Index}
\label{sec:posting}

\subsubsection{Posting Lists and Invariants}

Each term has a posting list ordered by document ID. A posting entry stores
$(\mathrm{docID},\mathrm{posStart},\mathrm{tf})$, where \code{posStart} is the
start of the entry's positions in a separate position array and \code{tf} is
the number of occurrences in that document. Positions are offsets in the
stop-word-filtered token stream, so a phrase query can verify adjacency without
returning to the original text.

The loader checks five structural invariants: document IDs are dense from 1 to
$N$; document IDs within one posting list are strictly increasing; the
\code{tf} of each entry equals its number of positions; the total number of
positions equals the total number of kept tokens; and the term table is in
strict lexicographic order. The same invariants are enforced by
\code{validateIndex()} after loading.

\subsubsection{Build Algorithm}

The builder scans the document stream and appends postings incrementally; at
the end, the dictionary is frozen into a sorted term table.

\begin{lstlisting}[style=cli]
for each document D in manifest order:
    pos = 0
    for each token t in tokenize(D):
        if stop_words.contains(t): continue
        term = stem(t) if stem_flag else t
        dict.upsert(term).append(D.id, pos)
        pos += 1
    index.docs.push({D.id, D.title, pos})
index.terms = dict.sorted_entries()
\end{lstlisting}

The formal configuration produces 14,937 terms, 117,704 posting entries,
556,866 positions, and a 3,685,655-byte v1 index. The file layout is shown in
Figure~\ref{fig:index-layout}.

% figures/index-layout.tex -- sequential index layout and posting structure.
% Pure TikZ diagram: field names and byte widths come from include/index.hpp
% and the v1 serializer in src/index.cpp (28 B header, 12+len doc entry,
% 28+len term entry, 8+4*tf posting entry). No measured values.
\begin{figure}[!htbp]
\centering
\begin{tikzpicture}[
  font=\footnotesize,
  >=Stealth,
  x=1cm, y=1cm,
  seglabel/.style={font=\footnotesize, text=mseblue, anchor=east},
  strip/.style={draw=black!50, rounded corners=1pt, line width=0.5pt,
                fill=black!4, anchor=west, align=left, inner sep=2.2pt,
                font=\scriptsize, text width=12.35cm},
  cell/.style={draw=black!55, rounded corners=1pt, line width=0.45pt,
               fill=white, align=center, inner sep=1.5pt,
               minimum height=0.52cm, font=\scriptsize},
  structbox/.style={draw=mseblue!70, rounded corners=3pt, line width=0.6pt,
                    fill=mseblue!4},
  title/.style={font=\small, text=black!80, anchor=west, align=left,
                text width=15.2cm},
  note/.style={font=\scriptsize, text=msegray, anchor=west, align=left},
]

% ======================================================= (a) sequential layout
\node[title] at (0.20,7.55)
  {(a) Sequential layout of \texttt{index.bin} (file order; lengths not to
   scale)};
\node[note, text width=15.2cm] at (0.20,7.10)
  {v2 (\texttt{-{}-compress}) changes only the posting segment: document IDs
   and positions become delta-coded varints, \texttt{version} is set to 2, and
   the other segments are unchanged.};

% --- header
\node[seglabel] at (2.42,6.30) {header (28 B)};
\node[strip] at (2.55,6.30)
  {Header, 28 B: magic \texttt{"MSE1"}, then five count and flag fields\\
   \texttt{version} u32 · \texttt{flags} u32 · \texttt{docCount} u32 ·\\
   \texttt{termCount} u32 · \texttt{tokenCount} u64};
% --- document table
\node[seglabel] at (2.42,5.35) {doc table};
\node[strip] at (2.55,5.35)
  {12 + len(title) B each, in increasing docID order\\
   \texttt{docID} u32 · \texttt{keptTokens} u32 · \texttt{titleLen} u32 ·
   \texttt{title} (variable bytes)};
% --- term table
\node[seglabel] at (2.42,4.40) {term table};
\node[strip] at (2.55,4.40)
  {28 + len(term) B each, in lexicographic term order
   (\texttt{offset} points to the posting block)\\
   \texttt{len} u32 · \texttt{term} · \texttt{df} u32 · \texttt{tfTotal} u64 ·
   \texttt{offset} u64 · \texttt{bytes} u32};
% --- postings
\node[seglabel] at (2.42,3.45) {postings};
\node[strip] at (2.55,3.45)
  {One block per term, 8·$\Sigma$df + 4·$\Sigma$tf B, no padding between
   blocks\\
   \texttt{df} × (\texttt{docID} u32 · \texttt{tf} u32) + $\Sigma$tf
   \texttt{positions} u32};

% ======================================================= (b) term -> block
\node[title] at (0.20,2.95)
  {(b) One term record $\rightarrow$ its posting block};

% --- term record: six field cells, byte width on the second line
\node[cell, minimum width=0.90cm] (len) at (0.85,2.20) {len\\[-1pt]{\scriptsize u32}};
\node[cell, minimum width=1.60cm] (term) at (2.15,2.20) {term\\[-1pt]{\scriptsize var. bytes}};
\node[cell, minimum width=0.90cm] (df) at (3.50,2.20) {df\\[-1pt]{\scriptsize u32}};
\node[cell, minimum width=1.30cm] (tftotal) at (4.70,2.20) {tfTotal\\[-1pt]{\scriptsize u64}};
\node[cell, minimum width=1.20cm] (offset) at (6.00,2.20) {offset\\[-1pt]{\scriptsize u64}};
\node[cell, minimum width=1.10cm] (bytes) at (7.20,2.20) {bytes\\[-1pt]{\scriptsize u32}};

\node[note, text width=7.90cm] at (0.35,1.48)
  {One term record = 28 + len(term) B; \texttt{offset} points to the first byte
   of that term's posting block.};

% --- posting block (two records shown; all fields are u32)
\draw[structbox] (9.55,0.38) rectangle (15.42,2.80);

\node[cell, minimum width=0.95cm] at (10.22,2.44) {docID};
\node[cell, minimum width=0.55cm] at (11.18,2.44) {tf};
\node[cell, minimum width=0.72cm] at (12.00,2.44) {pos$_0$};
\node[cell, minimum width=0.72cm] at (12.85,2.44) {pos$_1$};
\node[cell, minimum width=0.55cm] at (13.60,2.44) {$\cdots$};
\node[cell, minimum width=1.35cm] at (14.68,2.44) {pos$_{\mathrm{tf}-1}$};
\node[font=\scriptsize, text=msegray] at (12.54,2.02) {$\vdots$};
\node[cell, minimum width=0.95cm] at (10.22,1.60) {docID};
\node[cell, minimum width=0.55cm] at (11.18,1.60) {tf};
\node[cell, minimum width=0.72cm] at (12.00,1.60) {pos$_0$};
\node[cell, minimum width=1.15cm] at (13.30,1.60) {$\cdots$};
\node[cell, minimum width=1.35cm] at (14.68,1.60) {pos$_{\mathrm{tf}-1}$};

\node[note, anchor=north west, text width=5.40cm] at (9.76,1.24)
  {Each entry is 8 + 4·tf B: \texttt{docID}, \texttt{tf},\\
   then tf ascending u32 positions; df entries, packed.};

% --- term record -> posting block
\draw[->, line width=0.6pt, draw=black!65]
  (bytes.east) -- (9.55,2.20);
\node[note, anchor=south, inner sep=1pt] at (8.55,2.26) {\texttt{seek(offset)}};

\end{tikzpicture}
\caption{Index layout: \texttt{index.bin} stores the header, document table,
term table, and posting segment in order. A term record points to its posting
block with \texttt{offset} and \texttt{bytes}. Within a block, each posting
stores the document ID and term frequency followed by tf ascending positions,
so a phrase query only needs to compare adjacent positions.}
\label{fig:index-layout}
\end{figure}


\subsubsection{Persistence Formats v1 and v2}
\label{sec:persistence}

Persistence is part of the index specification, not an afterthought. Format v1
writes the header, document table, term table, and posting blocks in order. All
offsets are computed before writing, and the reader parses the same layout and
replays the invariant checks. The same index object therefore always produces
the same byte sequence.

Format v2 changes only the posting region: document IDs are stored as deltas,
positions as deltas, and frequencies and deltas use a 1--4 byte byte-aligned
variable-length code. The header, document table, and term table are unchanged,
except that \code{version} is set to 2. On the formal index, v2 reduces
3,685,655 bytes to 1,725,694 bytes, or 46.82\% of v1, with identical query
output.

The variable-length code uses the top two bits of the first byte for the total
length and stores 6, 14, 22, or 30 payload bits in 1, 2, 3, or 4 bytes,
respectively. Values are little-endian. A document-ID delta is at least 1; a
position delta is at least 0.

\subsection{Dictionary Structures}
\label{sec:dict}

The build-time dictionary exposes three operations: \code{upsert},
\code{find}, and \code{sortedEntries}. The primary implementation is a
self-implemented splay tree; a hash table and \code{std::map} are alternative
implementations used for comparison.

\subsubsection{Splay Tree}

The workload explains the choice. The index is insert-only, so the expensive
deletion cases of a red-black tree never occur. The token stream is
Zipf-distributed, and splaying an accessed node to the root keeps frequent terms
near the top. At the end of the build, an in-order traversal produces the sorted
term table directly, which is exactly the frozen dictionary the serializer
needs. The amortized cost of an operation is $O(\log T)$ for $T$ terms; the
worst case is $O(T)$ for a single operation.

\subsubsection{Hash Table and \texttt{std::map}}

The hash dictionary provides expected $O(1)$ lookup and no ordered traversal;
it is the speed baseline. The \code{std::map} dictionary is a red-black tree
with worst-case $O(\log T)$ lookup and stable behavior; it is the balanced-tree
control. Both implement the same three-operation interface, and all three
produce byte-identical index files. The comparison of their build costs is
reported in Section~\ref{sec:dict-analysis}.

\subsubsection{Comparison}

The splay tree pays for its ordered traversal with rotations and a higher
comparison count; the hash table is faster but cannot produce a sorted term
list without an extra sort; \code{std::map} sits between them with predictable
costs. None of these differences changes the persisted index, because the
dictionary is frozen before serialization.

\subsection{Disk-Based B+ Tree}
\label{sec:bptree-design}

For a vocabulary that does not fit in memory, the query side can use a disk B+
tree \emph{sidecar} that maps each term to its posting offset and block length.
The tree is a 4\,KiB page file: page 0 is the header, and each later page is a
leaf or an inner node. Leaves are linked by \code{nextLeaf} in key order. Keys
are variable-length strings; the fanout is determined by the byte budget.

The value of a leaf entry is a 64-bit integer that packs the posting block
offset and the block length:
\begin{equation}
  \mathrm{value}=(\mathrm{offset}\ll 21)\;|\;\mathrm{blockBytes}.
\end{equation}
The length field is 21 bits, which is ample for the block sizes in this
project.

The tree maintains five invariants, checked by \code{validate()}. The routing
interval is left-closed, $[\mathrm{key}_{i-1},\mathrm{key}_i)$, because a
separator key is a real copy of the smallest key in the right subtree. With
variable-length keys, a count-based split does not guarantee a 50\% byte fill;
the implementation splits at the byte midpoint of a legal window and uses the
provable lower bound \code{pageSize/2 - maxEntry} for the fill check. A lookup
reads the root-to-leaf path and then seeks to the posting block:

\begin{lstlisting}[style=cli]
page = readPage(root)
while page is an inner node:
    i = upper_bound(page.keys, key)   # left-closed intervals
    page = readPage(page.child[i])
binary search page.keys for key
if found: return page.value
else: return not_found
\end{lstlisting}

A lookup costs $O(h)$ page reads and $O(h\log F)$ comparisons for height $h$
and fanout $F$; it does not load the full term table. The sidecar is built by
bulk-loading the sorted terms and is rebuilt whenever the index is rebuilt.

\subsection{Block-Based External Construction}
\label{sec:spimi}

The in-memory builder keeps the whole dictionary and all postings in memory.
For a large corpus this is the wrong bound, so the project also implements a
SPIMI-style block builder. Documents are streamed one at a time; the buffer is
spilled to a sorted temporary \emph{run} when it reaches a term budget
(\code{--block-terms}) or an estimated byte budget (\code{--block-bytes}). Runs
are then merged in passes with at most \code{--merge-fan-in} inputs per pass.
The term metadata and postings are streamed to temporary files and assembled
into the final index, so the process does not rebuild the full index or all
compressed blocks in memory.

\begin{lstlisting}[style=cli]
for each document D in manifest order:
    add tokens of D to buffer
    if buffer.terms >= term_budget or buffer.bytes >= byte_budget:
        sort buffer; write run; clear buffer
if buffer is non-empty: sort buffer; write run
merge runs in passes of at most fan_in runs
\end{lstlisting}

The memory bound is the block budget plus document metadata, run paths, file
buffers, and one document. The budget is checked at document boundaries, so a
single document can exceed it. Because every run boundary falls between
documents, document IDs increase across runs, and the postings of the same term
can be concatenated during the merge without reordering.

On the Shakespeare corpus with a budget of 4,096, the builder produces 17 runs,
peaks at 5,129 live terms, writes 4,166,556 bytes of temporary run data, and
produces an index byte-identical to the in-memory build. The merge cost is
$O(T\cdot R)$ with a linear scan of all run cursors, or $O(T\log R)$ with a heap;
the implementation uses the simpler scan and reports the cost in
Section~\ref{sec:compression}.

\subsection{Query Processing and Threshold Design}
\label{sec:query-design}

\subsubsection{Normalization Pipeline}

The query text is tokenized, filtered, and stemmed by the same functions used at
build time. The index header flags decide whether stop-word filtering and
stemming are applied. Every term that is filtered, missing from the dictionary,
or dropped by a threshold is reported in an explanatory message, so an empty
result is not silent.

\subsubsection{Word, Phrase, and Boolean Queries}

A word query must normalize to exactly one term. A multi-word query uses
\code{--words}; a phrase query uses \code{--phrase}. For a phrase, the engine
starts from the term with the smallest document frequency, takes its documents
as candidates, and verifies that the remaining terms occur at
\code{start+1}, \code{start+2}, and so on. Positions are sorted, so each check
is a binary search.

\begin{lstlisting}[style=cli]
phrase_query(terms):
    choose t0 with the smallest df
    for each document d in postings(t0):
        for each start position p of t0 in d:
            if for every i in 1..k-1: p+i is in positions(t_i, d):
                emit d; break
\end{lstlisting}

A Boolean query requires all searchable terms by default. The score vector is
\code{(matched, tfSum, docID)}, sorted by descending matched count, descending
term-frequency sum, and ascending document ID, which makes the result order
deterministic. \code{--min-match k} requires at least $k$ matching terms; if
$k$ exceeds the number of searchable terms, it is clamped to that number.

\subsubsection{Thresholds}

All four thresholds act at query time; the index is built once. Their positions
and defaults are listed in Table~\ref{tab:threshold-params}.

\begin{table}[!ht]
  \centering
  \caption{Query thresholds: where they act and their defaults.}
  \label{tab:threshold-params}
  \small
  \begin{tabularx}{\linewidth}{l l l X}
    \toprule
    Option & Default & Stage & Meaning \\
    \midrule
    \code{--max-df} & 1.0 & query-term filter & drop a term whose df ratio exceeds the threshold \\
    \code{--min-df} & 1 & query-term filter & drop a term whose df is below the threshold \\
    \code{--min-match} & 0 & Boolean condition & require at least $k$ matching terms; 0 means all \\
    \code{--top-x} & 0 & result truncation & keep only the first $x$ documents after ranking \\
    \bottomrule
  \end{tabularx}
\end{table}

Phrase queries deliberately bypass \code{--max-df} and \code{--min-df}: a phrase
must match as a whole, and dropping any term changes its meaning. A query whose
terms are all stop words always returns nothing, because stop words never enter
the index; the program distinguishes ``stop word'' from ``not in the
dictionary'' in its explanation.

\subsection{Summary of Main Data Structures}
\label{sec:data-structures}

Table~\ref{tab:data-structures} summarizes the structures that carry the
invariants and costs discussed above. The table is a specification reference;
measured numbers appear in Chapter~\ref{sec:results}.

\begin{table}[!htbp]
  \centering
  \caption{Main data structures, invariants, and operation costs.}
  \label{tab:data-structures}
  \footnotesize
  \renewcommand{\arraystretch}{1.05}
  \begin{tabularx}{\linewidth}{@{}l X X X@{}}
    \toprule
    Structure & Key fields / layout & Invariants & Operations and cost \\
    \midrule
    Posting entry &
      \code{docID} u32, \code{posStart} u32, \code{tf} u32 &
      document IDs strictly increasing; \code{tf} equals the number of
      positions &
      append $O(1)$; phrase check $O(\log \mathrm{tf})$ per position \\
    Document table &
      \code{docID}, \code{keptTokens}, \code{title} &
      IDs dense from 1 to $N$; \code{keptTokens} sums to the total positions &
      lookup by ID $O(1)$ \\
    Splay dictionary &
      term, posting block, left/right child &
      binary-search-tree order; splay moves the accessed node to the root &
      \code{upsert}/\code{find} $O(\log T)$ amortized, $O(T)$ worst;
      \code{sortedEntries} $O(T)$ \\
    Hash dictionary &
      term, posting block, bucket array &
      load factor bounded; rehash preserves entries &
      expected $O(1)$ lookup; no ordered traversal \\
    B+ tree node &
      leaf flag, \code{keyCount}, \code{nextLeaf}, keys, values/children &
      left-closed routing intervals; byte fill lower bound; leaf chain in key
      order &
      lookup $O(h)$ page reads and $O(h\log F)$ comparisons \\
    Index header &
      \code{magic}, \code{version}, \code{flags}, \code{docCount},
      \code{termCount}, \code{tokenCount} &
      flags replay the build-time normalization; counts match the tables &
      read once $O(1)$ \\
    Block run &
      sorted sequence of (term, partial postings) &
      runs are disjoint by document range; document IDs increase across runs &
      merge $O(T\log R)$ with a heap, $O(TR)$ with a linear scan \\
    \bottomrule
  \end{tabularx}
\end{table}
